\documentclass[11pt]{article}
\usepackage[a4paper, top=18mm, bottom=18mm, left=20mm, right=20mm]{geometry}
\usepackage[T2A]{fontenc}
\usepackage[utf8]{inputenc}
\usepackage[ukrainian]{babel}
\usepackage{graphicx}
\usepackage[export]{adjustbox}
\usepackage{amsmath}
\usepackage{amssymb}
\graphicspath{{/app/Quiz/Scripts/2023/ProgAll/15BinTree/}}
\setlength{\parindent}{0pt}
\setlength{\parskip}{6pt}
\emergencystretch=3em
\pagestyle{empty}
\begin{document}
%begin
{\centering\Large\bfseries Контрольна робота\par}
\medskip
\noindent\rule{\linewidth}{0.5pt}
\smallskip
\noindent\textbf{Варіант \No\,1}\hfill \textit{ПІБ:}\,\underline{\hspace{60mm}}
\vspace{4mm}

%beginitem
1. %goodpath:Scripts/2018/Mod2018_1/03-good.txt
\textbf{Відмітити все, що є коректним оператором С++:}
( 1 ) 'a'=3;

( 2 ) cout>>3;

( 3 ) bool r=('N'>'5');

( 4 ) x+=y;
\vspace{5mm}
%enditem
%beginitem
2. \textbf{Що буде виведено за виконання фрагменту коду?}

%code
9**-1.0*True
%code
\vspace{5mm}
%enditem
%beginitem
3. \textbf{Що буде виведено за виконання фрагменту коду?}

%code
print(6 != 4 >= 6.0 == True)
%code
\vspace{5mm}
%enditem
%beginitem
4. %goodpath:Scripts/2019/Mod2019_1/Lex/01-good.txt
\textbf{Відмітити все, що є однією правильною лексемою:}
( 1 ) -3J

( 2 ) "a"a"

( 3 ) 0.5

( 4 ) 0b001
\vspace{5mm}
%enditem
%beginitem
5. \textbf{Що буде виведено за виконання фрагменту коду?}

a,b,c=4,2,1
def h(a):
global c    a=3
    b*=4
    c=5
    return a+b+c

a,b,c=5,6,0
print(h(b),a,b,c)
\vspace{5mm}
%enditem
%beginitem
6. \textbf{Що буде виведено за виконання фрагменту коду?}

%code
cout << (9 % 9 % 9 % 9);
%code
\vspace{5mm}
%enditem
%beginitem
7. 

\begin{center}
	\adjustimage{max size={0.8\linewidth}{0.8\paperheight}}
	{../15BinTree/trees_exp/133.png}
\end{center}
Указати послідовність міток вершин дерева, зображеного на рис., що утвориться в результаті обходу дерева в прямому порядку. Мітки записувати через пробіл.\par Алгоритм починає роботу з кореня (на рис. це верхня вершина).
%answer:133 прямому
\vspace{5mm}
%enditem
%beginitem
8. %goodpath:Scripts/2019/Mod2019_1/Lex/03-good.txt
\textbf{Відмітити все, що є коректним оператором С++:}
( 1 ) (2+2j)//2

( 2 ) if a<b: a=0 else: b=0

( 3 ) x=y=z

( 4 ) x==y==z
\vspace{5mm}
%enditem
%beginitem
9. \textbf{Що буде виведено за виконання фрагменту коду?}
%code
a, b, c = 9, 8, 7
def h(a, b=6, c=6):
    print(a, b, c, end=" ")

h(a=4, 2, c=1)
print(a, b, c)
%code
\vspace{5mm}
%enditem
%beginitem
10. \textbf{Що буде виведено за виконання фрагменту коду?}

print('R', end=' ')
e=[10,11,12,13,14]
e.extend([1,2])
print(e)

\vspace{5mm}
%enditem










%end
\newpage

\newpage

%begin
{\centering\Large\bfseries Контрольна робота\par}
\medskip
\noindent\rule{\linewidth}{0.5pt}
\smallskip
\noindent\textbf{Варіант \No\,2}\hfill \textit{ПІБ:}\,\underline{\hspace{60mm}}
\vspace{4mm}

%beginitem
1. \textbf{Що буде виведено за виконання фрагменту коду?}

a,b,c=4,2,1
def h(a):
global c    a=3
    b*=4
    c=5
    return a+b+c

a,b,c=5,6,0
print(h(b),a,b,c)
\vspace{5mm}
%enditem
%beginitem
2. \textbf{Що буде виведено за виконання фрагменту коду?}

%code
print('R', end=' ')
e = ('a','b','c','d','e',0,1,2,3)
print(e[-8])
%code

\vspace{5mm}
%enditem
%beginitem
3. %goodpath:Scripts/2018/Mod2018_1/03-good.txt
\textbf{Відмітити все, що є коректним оператором С++:}
( 1 ) 'a'=3;

( 2 ) cout>>3;

( 3 ) bool r=('N'>'5');

( 4 ) x+=y;
\vspace{5mm}
%enditem
%beginitem
4. \textbf{Що буде виведено за виконання фрагменту коду?}

%code
cout << (!6==4 || !6.0> true || 4.0<=3);
%code
\vspace{5mm}
%enditem
%beginitem
5. \textbf{Що буде виведено за виконання фрагменту коду?}
%code

a,b,c=9,8,7
def h(a,b=6,c=6):
    print(a,b,c,end="")

h(a=4,2,c=1)
print(a,b,c)

%code
\vspace{5mm}
%enditem
%beginitem
6. \textbf{Що буде виведено за виконання фрагменту коду?}

x,y,z=4,2,1
y,z,y,z=9,z,y,x
print(x,y,z)
\vspace{5mm}
%enditem
%beginitem
7. %goodpath:Scripts/2019/Mod2019_1/Lex/03-good.txt
\textbf{Відмітити все, що є коректним оператором С++:}
( 1 ) (2+2j)//2

( 2 ) if a<b: a=0 else: b=0

( 3 ) x=y=z

( 4 ) x==y==z
\vspace{5mm}
%enditem
%beginitem
8. %goodpath:Scripts/2018/Mod2018_1/01-good.txt
\textbf{Відмітити все, що є однією правильною лексемою:}
( 1 ) x*y

( 2 ) (1)

( 3 ) False

( 4 ) For
\vspace{5mm}
%enditem
%beginitem
9. \textbf{Що буде виведено за виконання фрагменту коду?}

%code
print(9**-1.0*True)
%code
\vspace{5mm}
%enditem
%beginitem
10. \textbf{Що буде виведено за виконання фрагменту коду?}

%code
#include <iostream>
using namespace std;

int h(int d){
    int y = 50;
    if (d == 0) 
        return 4;
    if (d > -2)
         y = 2;
    else
         y = 1;
    return y;
}

int main(){
    cout << h(-4);
    return 0;
}
%code

\vspace{5mm}
%enditem










%end
\newpage

\newpage

%begin
{\centering\Large\bfseries Контрольна робота\par}
\medskip
\noindent\rule{\linewidth}{0.5pt}
\smallskip
\noindent\textbf{Варіант \No\,3}\hfill \textit{ПІБ:}\,\underline{\hspace{60mm}}
\vspace{4mm}

%beginitem
1. \textbf{Що буде виведено за виконання фрагменту коду?}

%code
cout << 9 % 9 % 9 % 9;
%code
\vspace{5mm}
%enditem
%beginitem
2. %goodpath:Scripts/2019/Mod2019_1/Lex/01-good.txt
\textbf{Відмітити все, що є однією правильною лексемою:}
( 1 ) -3J

( 2 ) "a"a"

( 3 ) 0.5

( 4 ) 0b001
\vspace{5mm}
%enditem
%beginitem
3. \textbf{Що буде виведено за виконання фрагменту коду?}
try:
    res = 9**-1.0*True
    if isinstance(res, bool):
        print(f'bool;{res}')
    elif isinstance(res, int):
        print(f'int;{res}')
    elif isinstance(res, float):
        print(f'float;{res:.2f}')
    else:
        print('???')
except:
    print('Z')
\vspace{5mm}
%enditem
%beginitem
4. \textbf{Що буде виведено за виконання фрагменту коду?}
try:
    res = 8/8%8*8
    if isinstance(res, bool):
        print(f'bool;{res}')
    elif isinstance(res, int):
        print(f'int;{res}')
    elif isinstance(res, float):
        print(f'float;{res:.2f}')
    else:
        print('???')
except:
    print('Z')
\vspace{5mm}
%enditem
%beginitem
5. 
Вкажіть значення та тип виразу:
"6.0"==True
\vspace{5mm}
%enditem
%beginitem
6. %goodpath:Scripts/2019/Mod2019_1/Lex/03-good.txt
\textbf{Відмітити все, що є коректним оператором С++:}
( 1 ) (2+2j)//2

( 2 ) if a<b: a=0 else: b=0

( 3 ) x=y=z

( 4 ) x==y==z
\vspace{5mm}
%enditem
%beginitem
7. \textbf{Що буде виведено за виконання фрагменту коду?}
try:
    for e in range(10, 10, -3):
        if e>=7:
            break
        print(e, end=';')
        if e>=7:
            break
    else:
        print(e,end=';')
    print(e)
except:
    print('Z')
\vspace{5mm}
%enditem
%beginitem
8. \textbf{Що буде виведено за виконання фрагменту коду?}
%code
#include <iostream>
using namespace std;

int f(int &x, int &y){
    x = 6;
    y+= 4;
    return x;
}

int main(){
    int a = 5, b = 9;
    a = f(a, a);
    cout << a << ":" << b <<':';
    {
        int a = 1, b = 8;
        cout << ((a>=5) && ((b+=1) >= 5))
             << a << ":" << b << ":";
    }
    {
        inta = 3;
        cout << a << ':';
    }
    cout << a << ":" << b << endl;
    return 0;
}
%code

\vspace{5mm}
%enditem
%beginitem
9. %goodpath:Scripts/2018/Mod2018_1/01-good.txt
\textbf{Відмітити все, що є однією правильною лексемою:}
( 1 ) x*y

( 2 ) (1)

( 3 ) False

( 4 ) For
\vspace{5mm}
%enditem
%beginitem
10. \textbf{Що буде виведено за виконання фрагменту коду?}

%code
print('R', end=' ')
m=4
while m<10:
    m+=1
print(m)
%code

\vspace{5mm}
%enditem










%end
\newpage

\newpage

%begin
{\centering\Large\bfseries Контрольна робота\par}
\medskip
\noindent\rule{\linewidth}{0.5pt}
\smallskip
\noindent\textbf{Варіант \No\,4}\hfill \textit{ПІБ:}\,\underline{\hspace{60mm}}
\vspace{4mm}

%beginitem
1. \textbf{Що буде виведено за виконання фрагменту коду?}
a = 4
b = 2
c = 1

def h(a):
global c    a = 3
    b *= 4
    c = 5
    return a + b + c

a = 5
b = 6
c = 0

try:
    print(h(b), a, b, c, sep=';')
except:
    print('Z', a, b, c, sep=';')
\vspace{5mm}
%enditem
%beginitem
2. 
Записати ВИРАЗ, значенням якого є множина, що складається з цілих чисел від 15 до 2005 (включно), що не діляться на 7.\vspace{5mm}
%enditem
%beginitem
3. %goodpath:Scripts/2019/Mod2019_1/Lex/03-good.txt
\textbf{Відмітити все, що є коректним оператором С++:}
( 1 ) (2+2j)//2

( 2 ) if a<b: a=0 else: b=0

( 3 ) x=y=z

( 4 ) x==y==z
\vspace{5mm}
%enditem
%beginitem
4. 
Записати ВИРАЗ, значенням якого є список, що складається з кубів цілих чисел від 15 до 2005 (включно), що не діляться на жодне з чисел 2, 3, записаних в оберненому порядку.\vspace{5mm}
%enditem
%beginitem
5. %goodpath:Scripts/2018/Mod2018_1/01-good.txt
\textbf{Відмітити все, що є однією правильною лексемою:}
( 1 ) x*y

( 2 ) (1)

( 3 ) False

( 4 ) For
\vspace{5mm}
%enditem
%beginitem
6. 
Написати програму, що запитує у користувача значення дійсних параметрів $x$ та $y$, обчислює для них значення виразу 
$$\frac{\log_{2} x + 76}{(\arccos x + 0.6) (y + 87)}$$ 
та повiдомляє введенi данi та результат обчислення.
 Оцінюється правильність та якість коду.

\vspace{5mm}
%enditem
%beginitem
7. %goodpath:Scripts/2019/Mod2019_1/Lex/01-good.txt
\textbf{Відмітити все, що є однією правильною лексемою:}
( 1 ) -3J

( 2 ) "a"a"

( 3 ) 0.5

( 4 ) 0b001
\vspace{5mm}
%enditem
%beginitem
8. \textbf{Що буде виведено за виконання фрагменту коду?}

print('R', end=' ')
for e in range(10,10,-3):
    if e>=7:
        break
        print(e, end=' ')
    if e>=7:
        break
    else:
        print(e, end=' ')
    print(e)
\vspace{5mm}
%enditem
%beginitem
9. \textbf{Що буде виведено за виконання фрагменту коду?}
e=[10,11,12,13,14]; e.extend([1,2]); print(e)\vspace{5mm}
%enditem
%beginitem
10. 

\begin{center}
	\adjustimage{max size={0.8\linewidth}{0.8\paperheight}}
	{../15BinTree/trees_exp/133.png}
\end{center}
Указати послідовність міток вершин дерева, зображеного на рис., що утвориться в результаті обходу дерева в прямому порядку. Мітки записувати через пробіл.\par Алгоритм починає роботу з кореня (на рис. це верхня вершина).
%answer:133 прямому
\vspace{5mm}
%enditem










%end
\newpage
\end{document}


